\chapter{Giới thiệu và quy trình phát triển}\label{ch:gioithieu}
\cref{tab:doichieu-rutgon} cho biết nhóm Task của Assignment 1 nằm ở chương nào; bản đầy đủ ở \cref{app:doichieu}.
\begin{table}[htbp]\centering
\caption{Đối chiếu rút gọn giữa Task của đề và chương báo cáo}\label{tab:doichieu-rutgon}
\begin{tabular}{@{}p{6.4cm}p{6.6cm}@{}}\toprule
\textbf{Nội dung của đề} & \textbf{Chương} \\ \midrule
Tổ chức Scrum, Product Backlog, Sprint & \cref{ch:quanly} \\
Phân tích yêu cầu, use case & \cref{ch:yeucau} \\
Thiết kế kiến trúc, dữ liệu & \cref{ch:kientruc} \\
Task 1--10: triển khai & \cref{ch:trienkhai-se} \\
Kiểm thử và nghiệm thu & \cref{ch:kiemthu} \\
Bàn giao, Retrospective & \cref{ch:banggiao} \\ \bottomrule
\end{tabular}
\end{table}

\section{Bối cảnh, mục tiêu, phạm vi}\label{sec:se-boicanh}
Quán cà phê quy mô nhỏ cần để khách đặt món và thanh toán không dùng tiền mặt ngay trên điện thoại, tránh xếp hàng và nhầm đơn. BrewLite là ứng dụng web mobile-first cho phép xem thực đơn, chọn size và topping, đặt đơn, thanh toán mô phỏng và theo dõi trạng thái đơn. Bài tập lớn gồm 10 Task theo Agile Scrum, cùng sản phẩm tối thiểu (MVP) chạy được từ menu đến xác nhận thanh toán; mỗi Task có người phụ trách chính và người nghiệm thu khác.

Phạm vi chi tiết (trong, ngoài, không thực hiện) ở \cref{sec:se-phamvi}: bốn nghiệp vụ lõi của Task 10 (máy trạng thái đơn, thanh toán idempotent, kiểm soát tồn kho đồng thời, khuyến mãi và điểm thưởng) nằm trong phạm vi; cổng thanh toán thật và giao diện quản trị thì không.

\section{Mô hình quy trình áp dụng}\label{sec:se-quytrinh}
Nhóm áp dụng Agile Scrum~\cite{scrumguide2020} với ba Sprint ngắn (14--18/10, 19--22/10, 23--25/10): thời gian gấp, yêu cầu có thể điều chỉnh sau mỗi Sprint và cần sản phẩm chạy được sớm để kiểm thử. Quy trình gồm Product Backlog, Sprint Backlog, Increment và bốn sự kiện Sprint Planning, Daily Scrum, Sprint Review, Sprint Retrospective (xem \cref{ch:quanly}).

\section{Công nghệ và công cụ sử dụng}\label{sec:se-congnghe}
\begin{table}[htbp]\centering
\caption{Công nghệ và công cụ}\label{tab:congnghe}
\begin{tabular}{@{}p{3.2cm}p{9.6cm}@{}}\toprule
\textbf{Lớp} & \textbf{Công nghệ} \\ \midrule
Frontend & Next.js (App Router), mobile-first, manifest PWA cơ bản (chưa có service worker) \\
Backend & NestJS, class-validator, JWT, bcrypt \\
Cơ sở dữ liệu & PostgreSQL 16, Prisma 7 (adapter \texttt{pg}), migration và seed \\
Kiểm thử & Vitest (unit), Supertest trên PostgreSQL thật (tích hợp), Playwright (E2E) \\
Vận hành & Docker Compose, GitHub Actions, GHCR, Render \\ \bottomrule
\end{tabular}
\end{table}
Kho mã nguồn: \url{https://github.com/tthan-prog/BrewLite}. Web: \url{https://brewlite-web.onrender.com}. API: \url{https://brewlite-api.onrender.com/health}. Phiên bản bàn giao: tag \texttt{v0.1.2}.

\section{Cấu trúc báo cáo}\label{sec:se-cautruc}
\cref{ch:quanly} trình bày quản lý dự án; \cref{ch:yeucau} đặc tả yêu cầu; \cref{ch:kientruc} kiến trúc, dữ liệu và máy trạng thái; \cref{ch:trienkhai-se} triển khai từng Task; \cref{ch:kiemthu} kiểm thử; \cref{ch:banggiao} bàn giao và nhìn lại; \cref{ch:truyvet} ma trận truy vết; \cref{ch:ketluan-se} kết luận.
